\documentclass[12pt]{article}
\usepackage{booktabs}
\usepackage{array}
\begin{document}

\begin{table}[h!]
\centering
\caption{Panel A: Moody Leverage and Service Expenditure - Column (3) Debt/Rev}
\begin{tabular}{lc}
\toprule
 & Debt/Rev \\
 & (3) \\
\midrule
moody leverage x 1924--1926 & -0.09 \\
 & (0.05) \\
moody leverage x 1929--1933 & -0.18*** \\
 & (0.07) \\
moody leverage x 1934--1938 & -0.33*** \\
 & (0.08) \\
moody leverage x 1941--1943 & -0.20* \\
 & (0.12) \\
\midrule
City FE & \checkmark \\
Year FE & \checkmark \\
1930 Pop x Year & \checkmark \\
$\Delta$1920-30 Pop x Year & \checkmark \\
Revenue & \checkmark \\
Region x Year & \checkmark \\
\midrule
R-sq (within) & 0.72 \\
N & 3,748 \\
Mean(y) & 19.16 \\
SD(y) & 31.22 \\
\bottomrule
\end{tabular}
\begin{tablenotes}
\small
\item \textbf{Note:} This table presents results using the bond repayment shock variable following Equation 5.2. The leverage measure (specified in the header) is interacted with a shock variable to isolate the plausibly exogenous component of leverage (``moodyleverage''). The coefficients of interest capture the marginal effect of financial leverage conditional on having a high refinancing burden between 1930 and 1935. The outcome variable is real per-capita capital outlay, estimated using Poisson pseudo-maximum likelihood. Control variables include population size and growth, total non-debt revenue, and region-by-year fixed effects. Standard errors are shown in parentheses and are clustered at the city level.
\end{tablenotes}
\end{table}

\end{document}
\documentclass[12pt]{article}
\usepackage{booktabs}
\usepackage{array}
\begin{document}

\begin{table}[h!]
\centering
\caption{Panel A: Moody Leverage and Service Expenditure - Column (3) Debt/Rev}
\begin{tabular}{lc}
\toprule
 & Debt/Rev \\
 & (3) \\
\midrule
moody leverage x 1924--1926 & -0.09 \\
 & (0.05) \\
moody leverage x 1929--1933 & -0.18*** \\
 & (0.07) \\
moody leverage x 1934--1938 & -0.33*** \\
 & (0.08) \\
moody leverage x 1941--1943 & -0.20* \\
 & (0.12) \\
\midrule
City FE & \checkmark \\
Year FE & \checkmark \\
1930 Pop x Year & \checkmark \\
$\Delta$1920-30 Pop x Year & \checkmark \\
Revenue & \checkmark \\
Region x Year & \checkmark \\
\midrule
R-sq (within) & NA \\
N & 3,748 \\
Mean(y) & 19.16 \\
SD(y) & 31.22 \\
\bottomrule
\end{tabular}
\begin{tablenotes}
\small
\item \textbf{Note:} This table presents results using the bond repayment shock variable following Equation 5.2. The leverage measure (specified in the header) is interacted with a shock variable to isolate the plausibly exogenous component of leverage (``moodyleverage''). The coefficients of interest capture the marginal effect of financial leverage conditional on having a high refinancing burden between 1930 and 1935. The outcome variable is real per-capita capital outlay, estimated using Poisson pseudo-maximum likelihood. Control variables include population size and growth, total non-debt revenue, and region-by-year fixed effects. Standard errors are shown in parentheses and are clustered at the city level.
\end{tablenotes}
\end{table}

\end{document}
\documentclass[12pt]{article}
\usepackage{booktabs}
\usepackage{array}
\begin{document}

\begin{table}[h!]
\centering
\caption{Panel A: Moody Leverage and Service Expenditure - Column (3) Debt/Rev}
\begin{tabular}{lc}
\toprule
 & Debt/Rev \\
 & (3) \\
\midrule
moody leverage x 1924--1926 & -0.09 \\
 & (0.05) \\
moody leverage x 1929--1933 & -0.18*** \\
 & (0.07) \\
moody leverage x 1934--1938 & -0.33*** \\
 & (0.08) \\
moody leverage x 1941--1943 & -0.20* \\
 & (0.12) \\
\midrule
City FE & \checkmark \\
Year FE & \checkmark \\
1930 Pop x Year & \checkmark \\
$\Delta$1920-30 Pop x Year & \checkmark \\
Revenue & \checkmark \\
Region x Year & \checkmark \\
\midrule
R-sq (within) & NA \\
N & 3,748 \\
Mean(y) & 19.16 \\
SD(y) & 31.22 \\
\bottomrule
\end{tabular}
\begin{tablenotes}
\small
\item \textbf{Note:} This table presents results using the bond repayment shock variable following Equation 5.2. The leverage measure (specified in the header) is interacted with a shock variable to isolate the plausibly exogenous component of leverage (``moodyleverage''). The coefficients of interest capture the marginal effect of financial leverage conditional on having a high refinancing burden between 1930 and 1935. The outcome variable is real per-capita capital outlay, estimated using Poisson pseudo-maximum likelihood. Control variables include population size and growth, total non-debt revenue, and region-by-year fixed effects. Standard errors are shown in parentheses and are clustered at the city level.
\end{tablenotes}
\end{table}

\end{document}
\documentclass[12pt]{article}
\usepackage{booktabs}
\usepackage{array}
\begin{document}

\begin{table}[h!]
\centering
\caption{Panel A: Moody Leverage and Service Expenditure - Column (3) Debt/Rev}
\begin{tabular}{lc}
\toprule
 & Debt/Rev \\
 & (3) \\
\midrule
moody leverage x 1924--1926 & -0.09 \\
 & (0.05) \\
moody leverage x 1929--1933 & -0.18*** \\
 & (0.07) \\
moody leverage x 1934--1938 & -0.33*** \\
 & (0.08) \\
moody leverage x 1941--1943 & -0.20* \\
 & (0.12) \\
\midrule
City FE & \checkmark \\
Year FE & \checkmark \\
1930 Pop x Year & \checkmark \\
$\Delta$1920-30 Pop x Year & \checkmark \\
Revenue & \checkmark \\
Region x Year & \checkmark \\
\midrule
R-sq (within) & 0.57 \\
N & 3,748 \\
Mean(y) & 19.16 \\
SD(y) & 31.22 \\
\bottomrule
\end{tabular}
\begin{tablenotes}
\small
\item \textbf{Note:} This table presents results using the bond repayment shock variable following Equation 5.2. The leverage measure (specified in the header) is interacted with a shock variable to isolate the plausibly exogenous component of leverage (``moodyleverage''). The coefficients of interest capture the marginal effect of financial leverage conditional on having a high refinancing burden between 1930 and 1935. The outcome variable is real per-capita capital outlay, estimated using Poisson pseudo-maximum likelihood. Control variables include population size and growth, total non-debt revenue, and region-by-year fixed effects. Standard errors are shown in parentheses and are clustered at the city level.
\end{tablenotes}
\end{table}

\end{document}
\documentclass[12pt]{article}
\usepackage{booktabs}
\usepackage{array}
\begin{document}

\begin{table}[h!]
\centering
\caption{Panel A: Moody Leverage and Service Expenditure - Column (3) Debt/Rev}
\begin{tabular}{lc}
\toprule
 & Debt/Rev \\
 & (3) \\
\midrule
moody leverage x 1924--1926 & -0.09 \\
 & (0.05) \\
moody leverage x 1929--1933 & -0.18*** \\
 & (0.07) \\
moody leverage x 1934--1938 & -0.33*** \\
 & (0.08) \\
moody leverage x 1941--1943 & -0.20* \\
 & (0.12) \\
\midrule
City FE & \checkmark \\
Year FE & \checkmark \\
1930 Pop x Year & \checkmark \\
$\Delta$1920-30 Pop x Year & \checkmark \\
Revenue & \checkmark \\
Region x Year & \checkmark \\
\midrule
R-sq (within) & 0.56 \\
N & 3,748 \\
Mean(y) & 19.16 \\
SD(y) & 31.22 \\
\bottomrule
\end{tabular}
\begin{tablenotes}
\small
\item \textbf{Note:} This table presents results using the bond repayment shock variable following Equation 5.2. The leverage measure (specified in the header) is interacted with a shock variable to isolate the plausibly exogenous component of leverage (``moodyleverage''). The coefficients of interest capture the marginal effect of financial leverage conditional on having a high refinancing burden between 1930 and 1935. The outcome variable is real per-capita capital outlay, estimated using Poisson pseudo-maximum likelihood. Control variables include population size and growth, total non-debt revenue, and region-by-year fixed effects. Standard errors are shown in parentheses and are clustered at the city level.
\end{tablenotes}
\end{table}

\end{document}
\documentclass[12pt]{article}
\usepackage{booktabs}
\usepackage{array}
\begin{document}

\begin{table}[h!]
\centering
\caption{Panel A: Moody Leverage and Service Expenditure - Column (3) Debt/Rev}
\begin{tabular}{lc}
\toprule
 & Debt/Rev \\
 & (3) \\
\midrule
moody leverage x 1924--1926 & -0.09 \\
 & (0.05) \\
moody leverage x 1929--1933 & -0.18*** \\
 & (0.07) \\
moody leverage x 1934--1938 & -0.33*** \\
 & (0.08) \\
moody leverage x 1941--1943 & -0.20* \\
 & (0.12) \\
\midrule
City FE & \checkmark \\
Year FE & \checkmark \\
1930 Pop x Year & \checkmark \\
$\Delta$1920-30 Pop x Year & \checkmark \\
Revenue & \checkmark \\
Region x Year & \checkmark \\
\midrule
R-sq (within) & NA \\
N & 3,748 \\
Mean(y) & 19.16 \\
SD(y) & 31.22 \\
\bottomrule
\end{tabular}
\begin{tablenotes}
\small
\item \textbf{Note:} This table presents results using the bond repayment shock variable following Equation 5.2. The leverage measure (specified in the header) is interacted with a shock variable to isolate the plausibly exogenous component of leverage (``moodyleverage''). The coefficients of interest capture the marginal effect of financial leverage conditional on having a high refinancing burden between 1930 and 1935. The outcome variable is real per-capita capital outlay, estimated using Poisson pseudo-maximum likelihood. Control variables include population size and growth, total non-debt revenue, and region-by-year fixed effects. Standard errors are shown in parentheses and are clustered at the city level.
\end{tablenotes}
\end{table}

\end{document}
\documentclass[12pt]{article}
\usepackage{booktabs}
\usepackage{array}
\begin{document}

\begin{table}[h!]
\centering
\caption{Panel A: Moody Leverage and Service Expenditure - Column (3) Debt/Rev}
\begin{tabular}{lc}
\toprule
 & Debt/Rev \\
 & (3) \\
\midrule
moody leverage x 1924--1926 & -0.09 \\
 & (0.05) \\
moody leverage x 1929--1933 & -0.18*** \\
 & (0.07) \\
moody leverage x 1934--1938 & -0.33*** \\
 & (0.08) \\
moody leverage x 1941--1943 & -0.20* \\
 & (0.12) \\
\midrule
City FE & \checkmark \\
Year FE & \checkmark \\
1930 Pop x Year & \checkmark \\
$\Delta$1920-30 Pop x Year & \checkmark \\
Revenue & \checkmark \\
Region x Year & \checkmark \\
\midrule
R-sq (within) & NA \\
N & 3,748 \\
Mean(y) & 19.16 \\
SD(y) & 31.22 \\
\bottomrule
\end{tabular}
\begin{tablenotes}
\small
\item \textbf{Note:} This table presents results using the bond repayment shock variable following Equation 5.2. The leverage measure (specified in the header) is interacted with a shock variable to isolate the plausibly exogenous component of leverage (``moodyleverage''). The coefficients of interest capture the marginal effect of financial leverage conditional on having a high refinancing burden between 1930 and 1935. The outcome variable is real per-capita capital outlay, estimated using Poisson pseudo-maximum likelihood. Control variables include population size and growth, total non-debt revenue, and region-by-year fixed effects. Standard errors are shown in parentheses and are clustered at the city level.
\end{tablenotes}
\end{table}

\end{document}
\documentclass[12pt]{article}
\usepackage{booktabs}
\usepackage{array}
\begin{document}

\begin{table}[h!]
\centering
\caption{Panel A: Moody Leverage and Service Expenditure - Column (3) Debt/Rev}
\begin{tabular}{lc}
\toprule
 & Debt/Rev \\
 & (3) \\
\midrule
moody leverage x 1924--1926 & -0.09 \\
 & (0.05) \\
moody leverage x 1929--1933 & -0.18*** \\
 & (0.07) \\
moody leverage x 1934--1938 & -0.33*** \\
 & (0.08) \\
moody leverage x 1941--1943 & -0.20* \\
 & (0.12) \\
\midrule
City FE & \checkmark \\
Year FE & \checkmark \\
1930 Pop x Year & \checkmark \\
$\Delta$1920-30 Pop x Year & \checkmark \\
Revenue & \checkmark \\
Region x Year & \checkmark \\
\midrule
R-sq (within) & 0.09 \\
N & 3,748 \\
Mean(y) & 19.16 \\
SD(y) & 31.22 \\
\bottomrule
\end{tabular}
\begin{tablenotes}
\small
\item \textbf{Note:} This table presents results using the bond repayment shock variable following Equation 5.2. The leverage measure (specified in the header) is interacted with a shock variable to isolate the plausibly exogenous component of leverage (``moodyleverage''). The coefficients of interest capture the marginal effect of financial leverage conditional on having a high refinancing burden between 1930 and 1935. The outcome variable is real per-capita capital outlay, estimated using Poisson pseudo-maximum likelihood. Control variables include population size and growth, total non-debt revenue, and region-by-year fixed effects. Standard errors are shown in parentheses and are clustered at the city level.
\end{tablenotes}
\end{table}

\end{document}
\documentclass[12pt]{article}
\usepackage{booktabs}
\usepackage{array}
\begin{document}

\begin{table}[h!]
\centering
\caption{Panel A: Moody Leverage and Service Expenditure - Column (3) Debt/Rev}
\begin{tabular}{lc}
\toprule
 & Debt/Rev \\
 & (3) \\
\midrule
moody leverage x 1924--1926 & -0.09 \\
 & (0.05) \\
moody leverage x 1929--1933 & -0.18*** \\
 & (0.07) \\
moody leverage x 1934--1938 & -0.33*** \\
 & (0.08) \\
moody leverage x 1941--1943 & -0.20* \\
 & (0.12) \\
\midrule
City FE & \checkmark \\
Year FE & \checkmark \\
1930 Pop x Year & \checkmark \\
$\Delta$1920-30 Pop x Year & \checkmark \\
Revenue & \checkmark \\
Region x Year & \checkmark \\
\midrule
R-sq (within) & 0.08 \\
N & 3,748 \\
Mean(y) & 19.16 \\
SD(y) & 31.22 \\
\bottomrule
\end{tabular}
\begin{tablenotes}
\small
\item \textbf{Note:} This table presents results using the bond repayment shock variable following Equation 5.2. The leverage measure (specified in the header) is interacted with a shock variable to isolate the plausibly exogenous component of leverage (``moodyleverage''). The coefficients of interest capture the marginal effect of financial leverage conditional on having a high refinancing burden between 1930 and 1935. The outcome variable is real per-capita capital outlay, estimated using Poisson pseudo-maximum likelihood. Control variables include population size and growth, total non-debt revenue, and region-by-year fixed effects. Standard errors are shown in parentheses and are clustered at the city level.
\end{tablenotes}
\end{table}

\end{document}
